\documentclass[a4paper,11pt]{article}
\usepackage[T2A]{fontenc}
\usepackage[utf8]{inputenc}
\usepackage[russian]{babel}
\usepackage{amsmath,amssymb,mathtools}
\usepackage{helvet}
\renewcommand{\familydefault}{\sfdefault}
\usepackage{icomma}
\usepackage[margin=20mm,top=21mm,bottom=22mm]{geometry}
\usepackage{microtype}
\usepackage{tikz}
\usetikzlibrary{arrows.meta,positioning,shapes.geometric,calc}
\usepackage{tkz-euclide}
\usepackage{pgfplots}
\pgfplotsset{compat=1.18}
\usepackage{booktabs,tabularx,longtable,array}
\usepackage{enumitem}
\usepackage{xcolor}
\usepackage{fancyhdr}
\usepackage{setspace}
\usepackage{hyperref}
\definecolor{accent}{HTML}{008080}
\definecolor{darkaccent}{HTML}{004D4D}
\definecolor{textgray}{HTML}{333333}
\definecolor{muted}{HTML}{65717A}
\definecolor{burgundy}{HTML}{8C4C52}
\hypersetup{colorlinks=true,linkcolor=darkaccent,urlcolor=accent,pdfauthor={Лёвин Артём Александрович},pdftitle={Чек-лист. Повторение ОГЭ. София Кальней. 22.11.2026}}
\setstretch{1.22}
\setlength{\parindent}{0pt}
\setlength{\parskip}{6pt}
\setlength{\headheight}{14pt}
\setlist[enumerate]{leftmargin=*,itemsep=6pt,topsep=5pt}
\pagestyle{fancy}
\fancyhf{}
\fancyhead[L]{\small\color{muted}Математика \textbar\ Повторение ОГЭ}
\fancyhead[R]{\small\color{muted}София Кальней}
\fancyfoot[C]{\small\color{muted}\thepage}
\renewcommand{\headrulewidth}{0.25pt}
\renewcommand{\footrulewidth}{0pt}
\newcommand{\secname}[1]{\section*{\textcolor{darkaccent}{#1}}}
\newcommand{\checkline}[1]{\par\noindent\textcolor{accent}{\(\square\)}\hspace{0.5em}#1\par}
\newcommand{\important}[1]{\textcolor{darkaccent}{\textbf{#1}}}
\begin{document}
\thispagestyle{empty}
\begin{center}
\vspace*{22mm}
{\Large\bfseries\color{muted} МАТЕМАТИКА}\par
\vspace{22mm}
{\fontsize{29}{34}\selectfont\bfseries\color{darkaccent}Чек-лист}\par
\vspace{11mm}
{\Large\bfseries Повторение ОГЭ}\par
\vspace{5mm}
{\large Неравенства, корни, задачи на работу,\\ графики функций и геометрия}\par
\vspace{21mm}
{\Large София Кальней}\par
\vspace{8mm}
{\normalsize 22 ноября 2026 г.}\par
\vfill
{\normalsize Лёвин Артём Александрович}\par
{\small эксклюзивно для Софии Кальней}\par
\vspace{17mm}
\begin{tikzpicture}[x=1cm,y=1cm]
\draw[burgundy,line width=0.85pt] (-5.4,0) .. controls (-3.8,.63) and (-2.5,-.6) .. (0,0)
 .. controls (2.1,.55) and (3.8,-.48) .. (5.4,.06);
\end{tikzpicture}
\vspace{8mm}
\end{center}
\newpage

\secname{1. Квадратные неравенства и метод интервалов}
\checkline{Узнаю полный квадрат: \(a^2-2ab+b^2=(a-b)^2\).}
\checkline{Переношу все слагаемые в одну сторону и раскладываю выражение на множители.}
\checkline{Нахожу нули множителей, отмечаю их на числовой прямой и определяю знаки.}
\checkline{При строгом знаке \(<\) или \(>\) корни не включаю в ответ.}

\textbf{Пример с занятия.} Решите неравенство
\[
81-18x+x^2<\sqrt2\,(x-9).
\]
Заметим, что \(x^2-18x+81=(x-9)^2\). Тогда
\begin{align*}
(x-9)^2-\sqrt2(x-9)&<0,\\
(x-9)(x-9-\sqrt2)&<0.
\end{align*}
Нули множителей: \(x_1=9\), \(x_2=9+\sqrt2\). Поскольку коэффициент при \(x^2\) положителен, произведение отрицательно \emph{между} корнями.
\begin{center}
\begin{tikzpicture}[scale=1.0]
\draw[-{Stealth[length=2.7mm]},thick] (-4.6,0)--(4.7,0) node[right] {$x$};
\draw[burgundy,line width=2.2pt] (-.85,0)--(2.08,0);
\draw[fill=white,draw=burgundy,very thick] (-.85,0) circle (3pt);
\draw[fill=white,draw=burgundy,very thick] (2.08,0) circle (3pt);
\node[above] at (-3,0) {$+$};\node[above] at (.63,0) {$-$};\node[above] at (3.3,0) {$+$};
\node[below=6pt] at (-.85,0) {$9$};\node[below=6pt] at (2.08,0) {$9+\sqrt2$};
\end{tikzpicture}
\end{center}
\important{Ответ:} \(x\in(9;\,9+\sqrt2)\).

\textbf{Второй тип из занятия.} Решите
\[
4-4x+x^2\ge\sqrt3\,(2-x).
\]
\begin{align*}
(x-2)^2+\sqrt3(x-2)&\ge0,\\
(x-2)(x-2+\sqrt3)&\ge0.
\end{align*}
Корни: \(2-\sqrt3\) и \(2\). Произведение неотрицательно снаружи отрезка между корнями, включая концы.
\[
\boxed{x\in(-\infty;\,2-\sqrt3]\cup[2;\,+\infty).}
\]
\textcolor{burgundy}{\textbf{Важно.}} Проверяйте порядок корней: \(2-\sqrt3>0\), поскольку \(2>\sqrt3\), а \(4>3\).

\clearpage
\begin{center}
{\Large\bfseries\color{darkaccent}Алгоритм: метод интервалов}\par
\vspace{13mm}
\begin{tikzpicture}[node distance=11mm,>=Stealth, thick,
 start/.style={ellipse,draw=darkaccent,align=center,minimum width=48mm,minimum height=9mm,fill=white},
 task/.style={rounded corners=2pt,draw=accent,align=center,text width=112mm,inner sep=8pt,fill=white},
 arrow/.style={-{Stealth[length=3mm]},draw=darkaccent,line width=.8pt}]
\node[start] (begin) {Начало};
\node[task,below=of begin] (shift) {Перенесите все влево и приведите подобные слагаемые};
\node[task,below=of shift] (fact) {Разложите выражение на множители, если возможно};
\node[task,below=of fact] (zeros) {Приравняйте каждый множитель к нулю и упорядочьте найденные точки};
\node[task,below=of zeros] (signs) {Отметьте точки на числовой прямой, определите знаки на интервалах};
\node[task,below=of signs] (choose) {Выберите интервалы с нужным знаком; при \(\le\) или \(\ge\) включите допустимые нули};
\node[task,below=of choose] (write) {Запишите объединение подходящих промежутков и проверьте граничные точки};
\node[start,below=of write] (end) {Ответ};
\foreach \x/\y in {begin/shift,shift/fact,fact/zeros,zeros/signs,signs/choose,choose/write,write/end}
 \draw[arrow] (\x.south)--(\y.north);
\end{tikzpicture}
\end{center}
\clearpage

\secname{2. Сравнение выражений с корнями}
\checkline{Перед возведением неравенства в квадрат убеждаюсь, что обе части неотрицательны.}
\checkline{Если прибавляю или вычитаю число, знак неравенства сохраняется.}
\checkline{Если умножаю или делю на отрицательное число, знак меняется на противоположный.}
\checkline{Использую простые оценки квадратов, когда это сокращает решение.}

\textbf{Пример 1.} Докажите, что \(2-\sqrt3>0\). Это эквивалентно \(2>\sqrt3\). Обе части положительны, поэтому можно сравнить квадраты: \(4>3\). Утверждение доказано.

\textbf{Пример 2.} Сравните \(\sqrt2-1\) и \(\sqrt3-2\). Первое число положительно, поскольку \(\sqrt2>1\). Второе отрицательно, поскольку \(\sqrt3<2\). Следовательно,
\[
\boxed{\sqrt2-1>\sqrt3-2.}
\]
Это быстрее, чем двукратное возведение в квадрат.

\textbf{Пример 3.} Сравните \(\sqrt5\) и \(\sqrt7-1\). Прибавим \(1\) к обеим частям сравнения: достаточно установить, что \(\sqrt5+1>\sqrt7\). Обе величины положительны, и
\[
(\sqrt5+1)^2=6+2\sqrt5>7,
\]
поскольку \(2\sqrt5>1\). Значит, \(\sqrt5>\sqrt7-1\).

\textcolor{burgundy}{\textbf{Типичная ошибка.}} Из \(-3>-2\sqrt7\) при умножении на \(-1\) следует \(3<2\sqrt7\), а знак \(>\) превращается в \(<\).

\clearpage
\secname{3. Текстовые задачи на производительность}
\checkline{Обозначаю искомую скорость работы через \(x\) и указываю единицы измерения.}
\checkline{Использую формулу \(\text{время}=\dfrac{\text{объём работы}}{\text{производительность}}\).}
\checkline{Проверяю, кто работает быстрее, а кто выполняет больший объём работы.}
\checkline{Исключаю отрицательные скорости и значения, обращающие знаменатели в нуль.}

\textbf{Пример 1.} Первая труба пропускает на \(3\)~л/мин меньше второй. Резервуар объёмом \(260\)~л вторая заполняет на \(6\) минут быстрее. Найдите производительность второй трубы.
\begin{center}
\renewcommand{\arraystretch}{1.22}
\begin{tabularx}{\linewidth}{@{}l>{\centering\arraybackslash}X>{\centering\arraybackslash}X>{\centering\arraybackslash}X@{}}
\toprule
Труба & Производительность, л/мин & Работа, л & Время, мин\\\midrule
Первая & \(x-3\) & \(260\) & \(\dfrac{260}{x-3}\)\\
Вторая & \(x\) & \(260\) & \(\dfrac{260}{x}\)\\\bottomrule
\end{tabularx}
\end{center}
Здесь \(x>3\). Уравнение по времени:
\[
\frac{260}{x-3}-\frac{260}{x}=6
\quad\Longrightarrow\quad
x(x-3)=130.
\]
Получаем \(x^2-3x-130=0=(x-13)(x+10)\). С учётом \(x>3\): \(x=13\)~л/мин.

\textbf{Пример 2.} Первая труба наполняет \(770\)~л, вторая --- \(830\)~л за одно и то же время. Вторая пропускает на \(6\)~л/мин больше. Если первая пропускает \(x\)~л/мин, то
\[
\frac{770}{x}=\frac{830}{x+6}
\Longrightarrow\quad 770(x+6)=830x
\Longrightarrow\quad x=77.
\]
Значит, первая пропускает \(77\)~л/мин, вторая --- \(83\)~л/мин. Проверка: \(770/77=830/83=10\) минут.

\clearpage
\secname{4. Рациональная функция: выколотая точка}
\checkline{До сокращения дроби выписываю ограничения исходного знаменателя.}
\checkline{Строю график упрощённой функции и исключаю точки с запрещёнными абсциссами.}
\checkline{Для прямой \(y=kx\) ищу общие точки из уравнения \(f(x)=kx\).}
\checkline{При подсчёте решений учитываю и касание, и исключённую точку.}

\textbf{Пример с занятия.} Постройте график
\[
y=\frac{(x^2+1)(x-2)}{2-x}.
\]
Исходный знаменатель обращается в нуль при \(x=2\), поэтому \(x\ne2\). Поскольку \(2-x=-(x-2)\), для всех допустимых \(x\)
\[
y=-x^2-1,\qquad x\ne2.
\]
Вершина параболы: \((0;-1)\). Можно взять пары симметричных точек \((-1;-2)\), \((1;-2)\), \((-2;-5)\), \((2;-5)\). Последнюю точку \((2;-5)\) нужно \emph{выколоть}.
\begin{center}
\begin{tikzpicture}
\begin{axis}[
 width=0.78\linewidth,height=70mm,
 axis lines=middle,xmin=-3,xmax=3.7,ymin=-11,ymax=1,
 xtick={-3,-2,-1,0,1,2,3},ytick={-10,-8,-6,-4,-2,0},
 grid=both,minor tick num=1,grid style={gray!12},major grid style={gray!22},
 xlabel={$x$},ylabel={$y$},xlabel style={at={(axis description cs:1,.88)},anchor=south},
 ylabel style={at={(axis description cs:.48,1)},anchor=west},
 tick label style={font=\small},clip=false]
\addplot[accent,very thick,domain=-3:3.4,samples=200] {-x^2-1};
\addplot[only marks,mark=*,mark size=2pt,color=darkaccent] coordinates {(0,-1)};
\addplot[only marks,mark=o,mark size=3.6pt,mark options={line width=1.2pt,draw=burgundy,fill=white}] coordinates {(2,-5)};
\node[anchor=west,font=\small,burgundy] at (axis cs:2.10,-5.1) {$(2;-5)$};
\end{axis}
\end{tikzpicture}
\end{center}

\clearpage
\secname{5. Прямая \texorpdfstring{$y=kx$}{y=kx}: ровно одна общая точка}
Для пересечения с графиком предыдущей функции решаем
\[
-x^2-1=kx\quad\Longleftrightarrow\quad x^2+kx+1=0,\qquad x\ne2.
\]
\textbf{Случай 1 --- касание.} Дискриминант \(D=k^2-4\) равен нулю. Поэтому \(k=2\) или \(k=-2\). Соответствующие абсциссы \(x=-1\) и \(x=1\) допустимы.

\textbf{Случай 2 --- одна из двух точек исключена.} Условие прохождения прямой через выколотую точку \((2;-5)\):
\[
-5=2k\quad\Longrightarrow\quad k=-\frac52.
\]
Тогда \(x^2-\frac52x+1=(x-2)(x-\frac12)=0\). Корень \(x=2\) исключён, но \(x=\frac12\) остаётся. Ровно одна общая точка есть и в этом случае.
\[
\boxed{k\in\left\{-\frac52,-2,2\right\}.}
\]
\textcolor{burgundy}{\textbf{Важно.}} Одного критерия \(D=0\) недостаточно для графика с выколотой точкой.

\clearpage
\secname{6. Подобие треугольников}
\checkline{Ищу пары равных углов, в том числе при параллельных прямых.}
\checkline{Записываю соответствующие стороны в одном и том же порядке.}
\checkline{Если известна часть стороны, выражаю целую сторону как сумму частей.}

\textbf{Пример с занятия.} В треугольнике \(ABC\) точки \(M\in AB\), \(N\in BC\), причём \(MN\parallel AC\). Дано: \(MN=16\), \(AC=20\), \(NC=15\). Найдите \(BN\).
\begin{center}
\begin{tikzpicture}[scale=.87]
\tkzDefPoint(0,0){A}
\tkzDefPoint(7,0){B}
\tkzDefPoint(3,4){C}
\tkzDefPoint(1.4,0){M}
\tkzDefPoint(3.8,3.2){N}
\tkzDrawPolygon[thick,color=textgray](A,B,C)
\tkzDrawSegment[thick,color=accent](M,N)
\tkzDrawPoints[fill=darkaccent,size=3](A,B,C,M,N)
\tkzLabelPoints[below](A,B,M)
\tkzLabelPoints[above](C,N)
\tkzLabelSegment[above left,pos=.5](A,C){$20$}
\tkzLabelSegment[above right,pos=.53](M,N){$16$}
\tkzLabelSegment[above right,pos=.5](C,N){$15$}
\end{tikzpicture}
\end{center}
Из параллельности получаем \(\triangle BMN\sim\triangle BAC\). Пусть \(BN=x\); тогда \(BC=x+15\). Значит,
\[
\frac{BN}{BC}=\frac{MN}{AC},\qquad
\frac{x}{x+15}=\frac{16}{20}=\frac45.
\]
\(5x=4x+60\), откуда \(\boxed{BN=60}\).

\clearpage
\secname{7. Равнобедренная трапеция и биссектрисы}
\checkline{Помню: углы при каждом основании равнобедренной трапеции равны.}
\checkline{Использую параллельность оснований и свойства биссектрис.}
\checkline{Проверяю, можно ли получить прямоугольный треугольник.}

\textbf{Геометрическая задача с занятия.} В равнобедренной трапеции \(ABCD\) с большим основанием \(AD\) биссектрисы углов \(A\) и \(C\) пересекаются в точке \(F\). Биссектриса угла \(A\) пересекает \(CD\) в точке \(K\). Прямые \(AB\) и \(CF\) параллельны. Известно, что \(FK=4\sqrt3\). Найдите \(CF\).
\begin{center}
\begin{tikzpicture}[scale=.87]
\tkzDefPoint(0,0){A}
\tkzDefPoint(2,3.464){B}
\tkzDefPoint(4,3.464){C}
\tkzDefPoint(6,0){D}
\tkzDefPoint(3,1.732){F}
\tkzDefPoint(4.5,2.598){K}
\tkzDrawPolygon[thick,color=textgray](A,B,C,D)
\tkzDrawSegment[very thick,color=accent](A,K)
\tkzDrawSegment[very thick,color=burgundy](C,F)
\tkzDrawPoints[fill=darkaccent,size=3](A,B,C,D,F,K)
\tkzLabelPoints[below](A,D)
\tkzLabelPoints[above](B,C)
\tkzLabelPoint[below right](F){$F$}
\tkzLabelPoint[right](K){$K$}
\tkzMarkRightAngle[size=.20,color=burgundy](F,K,C)
\tkzLabelSegment[below right,pos=.42](F,K){$4\sqrt3$}
\end{tikzpicture}
\end{center}
\textbf{Как рассуждать.} Пусть \(\angle DAB=\alpha\). В равнобедренной трапеции \(\angle BCD=180^\circ-\alpha\). Так как \(CF\) --- биссектриса и \(AB\parallel CF\), получаем
\[
\alpha=\frac{180^\circ-\alpha}{2}
\quad\Longrightarrow\quad \alpha=60^\circ.
\]
Следовательно, \(AK\) образует с основанием \(AD\) угол \(30^\circ\). Луч \(DC\) образует с направлением \(AD\) угол \(120^\circ\). Значит, \(\angle FKC=90^\circ\). В треугольнике \(CFK\) угол при \(F\) равен \(30^\circ\), а \(CF\) --- гипотенуза. Тогда
\[
FK=CF\cos30^\circ=CF\,\frac{\sqrt3}{2},
\quad CF=\frac{2\cdot4\sqrt3}{\sqrt3}=\boxed{8}.
\]

\clearpage
\secname{Тренировочный блок: 10 заданий}
\textcolor{muted}{Путь А. Сложность возрастает постепенно. Решения записывайте последовательно; проверяйте ограничения и единицы измерения.}
\subsection*{Базовый уровень}
\begin{enumerate}
\item Решите неравенство: \((x-4)^2<3(x-4)\).
\item Решите неравенство: \(x^2-10x+25\ge\sqrt2(x-5)\).
\item Определите знак числа \(3-\sqrt8\). Обоснуйте ответ сравнением квадратов.
\item Сравните числа \(\sqrt3-1\) и \(\sqrt5-2\). Укажите, какое больше.
\end{enumerate}
\subsection*{Средний уровень}
\begin{enumerate}[start=5]
\item Первая труба пропускает на \(4\)~л/мин меньше второй. Резервуар объёмом \(320\)~л вторая наполняет на \(4\) минуты быстрее. Сколько литров в минуту пропускает вторая труба?
\item Первая труба наполняет ёмкость объёмом \(420\)~л, а вторая --- \(480\)~л за одинаковое время. Производительность второй трубы на \(3\)~л/мин больше. Найдите производительность первой трубы.
\item Постройте график \(y=\dfrac{(x^2+4)(x-3)}{3-x}\). Укажите координаты вершины и выколотой точки.
\item Для функции \(y=\dfrac{(x^2+1)(x-3)}{3-x}\) определите все \(k\), при которых прямая \(y=kx\) имеет с графиком функции ровно одну общую точку.
\end{enumerate}
\subsection*{Повышенный уровень}
\begin{enumerate}[start=9]
\item В треугольнике \(ABC\) отрезок \(MN\parallel AC\), где \(M\in AB\), \(N\in BC\). Найдите \(BN\), если \(MN=12\), \(AC=18\), \(NC=10\).
\item В равнобедренной трапеции \(ABCD\) с большим основанием \(AD\) биссектрисы углов \(A\) и \(C\) пересекаются в точке \(F\). Биссектриса угла \(A\) пересекает сторону \(CD\) в точке \(K\), причём \(AB\parallel CF\). Найдите \(CF\), если \(FK=3\sqrt3\).
\end{enumerate}

\clearpage
\secname{Ответы к тренировочному блоку}
\renewcommand{\arraystretch}{1.44}
\begin{center}
\begin{tabularx}{\linewidth}{@{}>{\centering\arraybackslash}p{12mm}X@{}}
\toprule
№ & Ответ\\
\midrule
1 & \(x\in(4;7)\)\\
2 & \(x\in(-\infty;5]\cup[5+\sqrt2;+\infty)\)\\
3 & \(3-\sqrt8>0\)\\
4 & \(\sqrt3-1>\sqrt5-2\)\\
5 & \(20\)~л/мин\\
6 & \(21\)~л/мин\\
7 & Вершина \((0;-4)\); выколотая точка \((3;-13)\)\\
8 & \(k\in\{-\frac{10}{3},-2,2\}\)\\
9 & \(BN=20\)\\
10 & \(CF=6\)\\
\bottomrule
\end{tabularx}
\end{center}
\vspace{16mm}
\secname{Итоговая памятка}
\begin{tabularx}{\linewidth}{@{}p{42mm}X@{}}
\toprule
Сигнал в задаче & Первый полезный шаг\\\midrule
Квадрат разности & Представьте трёхчлен как квадрат двучлена\\
Произведение \(<0\) & Найдите нули множителей и выберите отрицательные интервалы\\
Числа с корнями & Оцените знаки и только затем сравнивайте квадраты\\
Трубы и резервуары & Таблица: работа, производительность, время\\
Дробь с сокращением & Сначала ограничение исходного знаменателя\\
Одна общая точка & Проверьте касание и выколотые точки\\
Параллельный отрезок в треугольнике & Ищите подобные треугольники\\
Биссектрисы в трапеции & Найдите углы и признаки прямоугольного треугольника\\
\bottomrule
\end{tabularx}
\vfill
\begin{center}\small\textcolor{muted}{Лёвин Артём Александрович эксклюзивно для Софии Кальней}\end{center}
\end{document}
